% source: RKF/theorum/27_rsc_primitive_to_completion_hierarchy.md
\hypertarget{recognition-seam-primitive-to-completion-hierarchy}{%
\chapter{Recognition-Seam Primitive-to-Completion Hierarchy}\label{recognition-seam-primitive-to-completion-hierarchy}}

\hypertarget{why-this-note-exists}{%
\section{1. Why this note exists}\label{why-this-note-exists}}

The buried Recognition-Seam, morphic, UGD, Vedic/Kerala, Hilbert-completion and
spectral-Fredholm chapters contain useful results, but they live at different
levels of mathematical maturity. This note mines the reusable core without
promoting every proposed connection or commutator into a theorem.

The corrected hierarchy is:

\begin{verbatim}
Shunya / pre-distinction ground
-> cut
-> trans-cut morphism
-> stabilized recognition class of morphisms
-> clock-free transition differential
-> generated composition residue
-> recovered middle identity
-> recognition-Cauchy completion
-> cut-generated positive forms
-> target-faithful finite memory
-> spectral/Fredholm shadow.
\end{verbatim}

The Hilbert and Fredholm layers are therefore derived completion/representation
layers. They are not primitive.

\begin{center}\rule{0.5\linewidth}{0.5pt}\end{center}

\hypertarget{primitive-cut-recognition-datum}{%
\section{2. Primitive cut-recognition datum}\label{primitive-cut-recognition-datum}}

A primitive datum is

\[
\mathfrak G=(\mathfrak S,X,\mathfrak C_\Sigma,\mathcal R,\mathcal A),
\]

where:

\begin{verbatim}
S                 pre-distinction ground;
X                 appearance space generated by admitted cuts;
C_Sigma           first cut / reference-creation operation;
R                 recognition map;
A                 declared trans-cut arrows.
\end{verbatim}

The first distinction is

\[
\mathfrak C_\Sigma:
\mathfrak S\leadsto
(\mathcal D^-,\mathcal D^+,\Sigma).
\]

A trans-cut arrow is

\[
\tau_\Sigma:\mathcal D^-\longrightarrow\mathcal D^+.
\]

The cut is internal to the calculus. It is not an exceptional point at which
calculus is suspended.

\begin{center}\rule{0.5\linewidth}{0.5pt}\end{center}

\hypertarget{objects-are-stabilized-morphic-lineages}{%
\section{3. Objects are stabilized morphic lineages}\label{objects-are-stabilized-morphic-lineages}}

\hypertarget{definition-3.1-morphic-object}{%
\subsection{Definition 3.1 (Morphic object)}\label{definition-3.1-morphic-object}}

An object is a stabilized recognition class of composable arrows:

\[
\boxed{
O=
[\alpha_n\cdots\alpha_2\alpha_1]_{\mathcal R}.
}
\]

The identity of \(O\) is the minimal closed arrow preserving its recovered
recognition state. It is not defined as the absence of change.

This imports the strongest part of the morphic chapter:

\begin{verbatim}
form is a recognizable stabilization of flow;
flow is prior to the stabilized object;
silence/projection is a repair channel, not object annihilation.
\end{verbatim}

No specific Tamil linguistic commutator is consumed as a theorem here. Those
remain domain examples until represented by declared operators.

\begin{center}\rule{0.5\linewidth}{0.5pt}\end{center}

\hypertarget{clock-free-transition-differential}{%
\section{4. Clock-free transition differential}\label{clock-free-transition-differential}}

For a declared arrow

\[
\gamma:x\longrightarrow y,
\]

and a form assignment \(G_x\), define

\[
\boxed{
\Delta_\gamma G
=A_\gamma^*G_yA_\gamma-G_x.
}
\]

For composable arrows \(\gamma:x\to y\), \(\delta:y\to z\),

\[
\boxed{
\Delta_{\delta\gamma}G
=
\Delta_\gamma G
+A_\gamma^*(\Delta_\delta G)A_\gamma.
}
\]

This is the native chain rule. It compares declared transitions and needs no
external clock or infinitesimal quotient.

A generated ledger \(\Lambda_\gamma^G\) is admissible only if it is declared
before the evaluation packet and obeys the same composition law:

\[
\Lambda_{\delta\gamma}^G
=
\Lambda_\gamma^G
+A_\gamma^*\Lambda_\delta^GA_\gamma.
\]

The open residue is

\[
\operatorname{Res}_\gamma^G
=
\Delta_\gamma G-\Lambda_\gamma^G.
\]

A ledger selected after seeing the residue is retrospective cancellation and
is inadmissible.

\begin{center}\rule{0.5\linewidth}{0.5pt}\end{center}

\hypertarget{generated-cut-corner-residues}{%
\section{5. Generated cut-corner residues}\label{generated-cut-corner-residues}}

After a positive representation is derived, let \(P_x\) and \(Q_x=I-P_x\)
be the recognized and memory cuts. For an event transport \(U_\gamma\), define

\[
\mathsf R_\gamma=P_yU_\gamma P_x,
\qquad
\mathsf J_\gamma=P_yU_\gamma Q_x,
\]

\[
\mathsf C_\gamma=Q_yU_\gamma P_x,
\qquad
\mathsf M_\gamma=Q_yU_\gamma Q_x.
\]

Then

\[
\mathsf R_{\delta\gamma}
=
\mathsf R_\delta\mathsf R_\gamma
+
\mathsf J_\delta\mathsf C_\gamma,
\]

\[
\mathsf M_{\delta\gamma}
=
\mathsf C_\delta\mathsf J_\gamma
+
\mathsf M_\delta\mathsf M_\gamma.
\]

The cross terms are generated composition residues. They are created by the
cut-corner product and therefore cannot be deleted merely because endpoint
shadows agree.

\begin{center}\rule{0.5\linewidth}{0.5pt}\end{center}

\hypertarget{recovered-middle-identity}{%
\section{6. Recovered middle identity}\label{recovered-middle-identity}}

Shadow equality of the middle object does not license composition.

\hypertarget{definition-6.1-recovered-middle-identity}{%
\subsection{Definition 6.1 (Recovered middle identity)}\label{definition-6.1-recovered-middle-identity}}

Suppose

\[
A\xrightarrow{\gamma}B_1,
\qquad
B_2\xrightarrow{\delta}C.
\]

The two middle states are recovered-identical when:

\begin{verbatim}
shadow state agrees;
phase transport is aligned;
scale/support is equivalent;
cut-memory transport closes or is lawfully ledgered;
path holonomy closes or is lawfully ledgered.
\end{verbatim}

Write

\[
B_1=_{\RSC}B_2.
\]

\hypertarget{theorem-6.2-recognition-composable-middle}{%
\subsection{Theorem 6.2 (Recognition-composable middle)}\label{theorem-6.2-recognition-composable-middle}}

The composition \(\delta\gamma\) is admitted only after recovered middle
identity. Endpoint equality or nonzero scalar overlap alone is insufficient.

This theorem is the morphic safeguard needed by every refinement sequence. It
prevents one finite packet from lending its visible coordinates to another
packet while silently changing the source metric, cut memory or normalization.

\begin{center}\rule{0.5\linewidth}{0.5pt}\end{center}

\hypertarget{recognition-completion-is-stronger-than-shadow-completion}{%
\section{7. Recognition completion is stronger than shadow completion}\label{recognition-completion-is-stronger-than-shadow-completion}}

A Hilbert or Banach limit is only a shadow limit unless the cut channels also
close.

\hypertarget{definition-7.1-recognition-cauchy-net}{%
\subsection{Definition 7.1 (Recognition-Cauchy net)}\label{definition-7.1-recognition-cauchy-net}}

A directed family of transported packets

\[
\widetilde x_i=(x_i,\phi_i,\sigma_i,k_i,\gamma_i)
\]

is recognition-Cauchy when:

\[
\|x_j-x_i\|\longrightarrow0,
\]

and, after recovered middle identification,

\begin{verbatim}
phase mismatch -> 0 or is ledgered;
scale/support mismatch -> 0;
cut-memory mismatch -> 0 or is ledgered;
transition residue -> 0;
path holonomy -> 0 or is ledgered.
\end{verbatim}

The completion is formed only after quotienting the positive recognition form
by its null seam. The pre-Hilbert/Hilbert carrier is therefore derived from the
recognition form; it is not primitive.

\hypertarget{theorem-7.2-shadow-completion-is-not-recognition-completion}{%
\subsection{Theorem 7.2 (Shadow completion is not recognition completion)}\label{theorem-7.2-shadow-completion-is-not-recognition-completion}}

Norm convergence alone does not imply recognition convergence. A constant
shadow vector may carry drifting cut memory, unaligned phase or nonzero loop
holonomy. Therefore an infinite-dimensional proof must certify both the
shadow tail and the recognition tail.

\begin{center}\rule{0.5\linewidth}{0.5pt}\end{center}

\hypertarget{madhavasmriti-completion-rule}{%
\section{8. Madhava--Smriti completion rule}\label{madhavasmriti-completion-rule}}

The Vedic/Kerala bridge supplies a useful convergence grammar:

\[
\mathcal R(S_n+C_n)-B_\infty-\mathsf S_n=0.
\]

Here \(S_n\) is the finite recursion, \(C_n\) is a predeclared correction,
\(B_\infty\) is the recognized limit and \(\mathsf S_n\) is the lawful tail
memory.

\hypertarget{definition-8.1-smriti-complete-approximation}{%
\subsection{Definition 8.1 (Smriti-complete approximation)}\label{definition-8.1-smriti-complete-approximation}}

A finite approximation is Smriti-complete if

\[
\boxed{
B_\infty
=
\mathcal R(S_n+C_n)-\mathsf S_n
}
\]

with

\[
\|\mathsf S_n\|\longrightarrow0
\]

in the declared recognition norm, and with \(C_n\) generated independently of
the evaluation residue.

This is the lawful form of

\begin{verbatim}
limit = finite recursion + correction + declared tail memory.
\end{verbatim}

It will be used for operator tails, source-event tails, determinant series and
prime/Gamma truncations. A visually convergent finite sequence without a
Smriti tail is only a shadow approximation.

\begin{center}\rule{0.5\linewidth}{0.5pt}\end{center}

\hypertarget{what-is-adopted-now}{%
\section{9. What is adopted now}\label{what-is-adopted-now}}

The following buried results are adopted as foundational inputs:

\begin{verbatim}
cut and trans-cut morphism;
objects as stabilized morphic lineages;
clock-free transition differential;
generated composition residues;
recovered middle identity;
projection leakage as open residue;
recognition-Cauchy completion;
Madhava--Smriti tail completion;
shadow versus recognition spectral data.
\end{verbatim}

The following are retained as useful representation charts, not primitives:

\begin{verbatim}
UGD phase-scale-seam coordinates (phi,sigma,k);
R^2=-I circular rotation and K^2=I hyperbolic seam chart;
exponential change packets;
Crown tensor connection;
classical Hilbert/Fredholm realizations.
\end{verbatim}

The following are not promoted without additional representation theorems:

\begin{verbatim}
specific morphic commutator constants;
existence of every proposed Crown connection component;
recognition-index direct sums with undefined target groups;
curvature-flux identities without a declared connection and domain;
physical or linguistic interpretations of formal operator symbols.
\end{verbatim}

\begin{center}\rule{0.5\linewidth}{0.5pt}\end{center}

\hypertarget{consequence-for-the-rh-programme}{%
\section{10. Consequence for the RH programme}\label{consequence-for-the-rh-programme}}

The completed-Weil specialization must be built in this order:

\begin{verbatim}
native trans-cut event morphisms
-> stabilized completed-Weil event object
-> cut-generated recognized and memory forms
-> recovered identity across all finite refinements
-> Smriti-complete prime/Gamma/boundary tails
-> recognition-complete limit
-> target-faithful memory observer of rank at most five
-> finite seam matrix
-> classical explicit-formula recognition only at the final shadow interface.
\end{verbatim}

The immediate next theorem is the Recognition-Complete Finite-to-Infinite Cut
Theorem. It must show that finite cut packets, once transported through
recovered identities and supplied with vanishing Smriti tails, converge to one
cut-generated infinite object with no hidden adverse complement.

\begin{center}\rule{0.5\linewidth}{0.5pt}\end{center}

\hypertarget{claim-boundary}{%
\section{11. Claim boundary}\label{claim-boundary}}

\begin{verbatim}
primitive cut/trans-cut hierarchy                     ADOPTED
clock-free transition differential                    PROVED ALGEBRAICALLY
generated cut-corner composition residues              PROVED ALGEBRAICALLY
recovered middle identity requirement                  ADOPTED AS COMPOSITION GATE
recognition completion stronger than norm completion   PROVED BY EXACT COUNTEREXAMPLES
Madhava--Smriti completion grammar                     ADOPTED AS TAIL DISCIPLINE

specific Crown/morphic physical models                 NOT PROMOTED
RH event lift                                           NOT YET CONSTRUCTED
RH recognition-complete tail theorem                    NEXT
RH                                                      NOT CLAIMED
\end{verbatim}

\begin{flushright}\small\texttt{RKF/theorum/27\_rsc\_primitive\_to\_completion\_hierarchy.md}\end{flushright}
